\documentclass[11pt,a4paper]{article}
\input{../../sablony/pripravy/priprava.sty}
\usepackage{tikz}
\nastavPripravu{1. ročník}{Souhrnné opakování}{135}{Matematika v praxi: technika a život}
\begin{document}
\blok{Úvod a připomenutí}
Aplikační stanoviště napříč obory: elektrotechnika (ELE), informatika (ICT), mechanika (MEC), ekonomika (EKO), chemie (CHE) a každodenní život (ZIV). \textbf{Číselná data jsou modelová, pokud není řečeno jinak.} Převádíme slovní situace na známé rovnice, soustavy, nerovnice nebo kvadratické funkce; žádné nové učivo.
\blok{Společně řešené příklady}
\begin{prikladReseniBox}[title={ELE: servis napájecího obvodu}]
Ideální stejnosměrný zdroj má $U=24\,\mathrm V$. Dva sériově zapojené odpory mají hodnoty $3\,\Omega$ a $R\,\Omega$. Při proudu $I=2\,\mathrm A$ určete $R$.
\tcblower
\[
\begin{aligned}
I&=\frac{U}{3+R}\\
2&=\frac{24}{3+R}\\
6+2R&=24\\
R&=9\,\Omega
\end{aligned}
\]

\end{prikladReseniBox}
\begin{prikladReseniBox}[title={EKO: kdy se vyplatí tisk?}]
3D tiskárna má modelové fixní náklady $250$ Kč a $12$ Kč na výrobek. Prodejní cena jednoho výrobku je $32$ Kč. Kolik \emph{celých} kusů musíme prodat, aby zisk byl kladný?
\tcblower
\[
\begin{aligned}
32n&>250+12n\\
20n&>250\\
n&>12{,}5\\
n&\ge13\quad(n\in\mathbb N_0)
\end{aligned}
\]

\end{prikladReseniBox}
\begin{prikladReseniBox}[title={MEC: let modelového tělesa}]
Výška nad zemí je v metrech dána $h(t)=-5t^2+20t$ pro $0\le t\le4$ s časem $t$ v sekundách. Určete maximální výšku, čas dosažení maxima a dobu návratu na zem.
\tcblower
\[
\begin{aligned}
h(t)&=-5(t-2)^2+20\\
h_{\max}&=20\,\mathrm m\quad\text{při }t=2\,\mathrm s\\
h(t)=0&\ \Rightarrow\ -5t(t-4)=0\\
\text{návrat }&t=4\,\mathrm s
\end{aligned}
\]

\end{prikladReseniBox}
\begin{prikladReseniBox}[title={ICT: který algoritmus vítězí?}]
Pro modelové běhy platí $T_A(n)=8+0{,}4n$ ms a $T_B(n)=2+0{,}7n$ ms. Při jakém $n\ge0$ se časy rovnají a který algoritmus je rychlejší pro malé $n$?
\tcblower
\[
\begin{aligned}
8+0{,}4n&=2+0{,}7n\\
6&=0{,}3n\\
n&=20\\
n<20&:\ B\text{ rychlejší}\\
n>20&:\ A\text{ rychlejší}
\end{aligned}
\]

\end{prikladReseniBox}
\blok{Samostatné procvičování}
\begin{zadaniReseniBox}[title={CHE: teoretický výtěžek hašení vápna}]
Při reakci $\mathrm{CaO+H_2O\rightarrow Ca(OH)_2}$ předpokládejme přebytek vody. Kolik gramů hydroxidu vápenatého teoreticky vznikne ze $56$ g oxidu vápenatého, použijeme-li molární hmotnosti přibližně $56$ a $74\,\mathrm{g/mol}$?
\tcblower
\[
\begin{aligned}
n(\mathrm{CaO})&=\frac{56}{56}=1\,\mathrm{mol}\\
n(\mathrm{Ca(OH)_2})&=1\,\mathrm{mol}\\
m(\mathrm{Ca(OH)_2})&=1\cdot74=74\,\mathrm g
\end{aligned}
\]

\end{zadaniReseniBox}
\begin{zadaniReseniBox}[title={ZIV: tolerance teploty}]
Zařízení má pracovat při teplotě nejvýše o $1{,}5\,^{\circ}\mathrm C$ odchýlené od $40\,^{\circ}\mathrm C$. Vyjádřete přípustné teploty intervalem.
\tcblower
\[
\begin{aligned}
|T-40|&\le1{,}5\\
K&=\langle38{,}5;41{,}5\rangle\,^{\circ}\mathrm C
\end{aligned}
\]

\end{zadaniReseniBox}
\begin{zadaniReseniBox}[title={CHE: směs dvou roztoků}]
Chceme $2$ litry roztoku o koncentraci $6\%$ ze dvou roztoků o koncentracích $3\%$ a $9\%$. Kolik litrů od každého potřebujeme? Předpokládejme sčítání objemů.
\tcblower
\[
\begin{aligned}
x+y&=2\\
0{,}03x+0{,}09y&=0{,}06\cdot2\\
3x+9y&=12\\
3(2-y)+9y&=12\\
x&=1\,\mathrm l,\quad y=1\,\mathrm l
\end{aligned}
\]

\end{zadaniReseniBox}
\blok{Volitelná rezerva a výzva}
\poznamka{Výzva: model CPI/počet taktů procesoru; rozhodni, jaké další údaje by byly potřeba pro reálné měření rychlosti hradel.}
\blok{Domácí úkol}
\begin{zadaniReseniBox}[title={Domácí příprava}]
\begin{enumerate}[label=\alph*),itemsep=2pt]
\item \textbf{EKO:} Výrobek stojí $180$ Kč fixně a $9$ Kč za kus, prodává se za $24$ Kč. Kolik celých kusů musí být prodáno, aby zisk dosáhl alespoň $120$ Kč?
\item \textbf{ENE:} Baterie o energii $96\,\mathrm{Wh}$ napájí spotřebič se stálým odběrem $12\,\mathrm W$. Jak dlouho vydrží v ideálním modelu?
\item \textbf{LOG:} Deset nákladů o hmotnostech $4$ a $7$ kg váží dohromady $58$ kg. Kolik je kterých?
\item \textbf{GEO:} Obdélník má obvod $30$ m. Jaké rozměry maximalizují jeho obsah a jaké je maximum?
\item \textbf{ICT:} Modelově platí $2t_{\mathrm{NAND}}+t_{\mathrm{XOR}}=13$ ns a $t_{\mathrm{NAND}}+2t_{\mathrm{XOR}}=14$ ns. Určete obě zpoždění.
\item \textbf{CHE:} Připravte $3$ litry roztoku o koncentraci $4\%$ z roztoků $2\%$ a $8\%$. Předpokládejte aditivitu objemů.
\end{enumerate}
\tcblower
\begin{enumerate}[label=\alph*),itemsep=2pt]
\item
\[
\begin{aligned}
24n-(180+9n)&\ge120\\
15n&\ge300\\
n&\ge20
\end{aligned}
\]
\item
\[
\begin{aligned}
t&=\frac{96\,\mathrm{Wh}}{12\,\mathrm W}\\
t&=8\,\mathrm h
\end{aligned}
\]
\item
\[
\begin{aligned}
x+y&=10\\
4x+7y&=58\\
4(10-y)+7y&=58\\
y&=6,\quad x=4
\end{aligned}
\]
\item
\[
\begin{aligned}
x+y&=15\\
S(x)&=x(15-x)\\
&=-(x-7{,}5)^2+56{,}25\\
\max S&=56{,}25\,\mathrm{m^2}\\
x&=y=7{,}5\,\mathrm m
\end{aligned}
\]
\item
\[
\begin{aligned}
2a+b&=13,\quad a+2b=14\\
3b&=15\ \Rightarrow\ b=5\,\mathrm{ns}\\
a&=4\,\mathrm{ns}
\end{aligned}
\]
\item
\[
\begin{aligned}
x+y&=3\\
0{,}02x+0{,}08y&=0{,}12\\
x&=2\,\mathrm l,\quad y=1\,\mathrm l
\end{aligned}
\]
\end{enumerate}
\end{zadaniReseniBox}
\poznamka{Čas: výběr dvou společných aplikací 16 min, stanoviště 20 min, rozbor a porovnání modelů 9 min. Ne všechny úlohy musí být vyřešeny v hodině. U chemie jde o teoretický výpočet, nikoli doporučení k provedení reakce; zaokrouhlené molární hmotnosti jsou pouze výukové.}
\end{document}
